% Colón-González et al. (2021): ვრცელი ქართული შეჯამება (სასწავლო მასალა, მხოლოდ სერგისთვის)
\documentclass[11pt,a4paper]{article}
\usepackage{fontspec}
\usepackage[a4paper,margin=2.3cm]{geometry}
\usepackage{booktabs}
\usepackage{longtable}
\usepackage{array}
\usepackage{xcolor}
\usepackage{enumitem}
\usepackage[hidelinks]{hyperref}
\usepackage{tcolorbox}

\setmainfont{Sylfaen}
\setmonofont{Consolas}

ewfontfamily\symfont{Segoe UI Symbol}

ewcommand{\ok}{{\symfont ✓}}
ewcommand{
o}{{\symfont ✗}}
ewcommand{\warn}{{\symfont ⚠}}

\newcommand{\q}[1]{,,#1``}
\newtcolorbox{mnishvn}{colback=orange!8,colframe=orange!70!black,title=რატომ არის ეს შენთვის მნიშვნელოვანი,fonttitle=\bfseries}
\newtcolorbox{citata}{colback=blue!4,colframe=blue!40!black,boxrule=0.4pt}
\newtcolorbox{ganmarteba}{colback=green!5,colframe=green!40!black,title=ტერმინის ახსნა,fonttitle=\bfseries}

\title{Colón-González et al.\ (2021)\\[0.3em]
\large ვრცელი შეჯამება: რაზეა სტატია და რა არის მნიშვნელოვანი შენი ნაშრომისთვის}
\author{სასწავლო მასალა სერგისთვის}
\date{2026-11-19}

\begin{document}
\maketitle

\noindent\textbf{სტატია:} Colón-González, F.\,J., Soares Bastos, L., Hofmann, B., \ldots, Brady, O.\,J., Lowe, R.\ (2021).
\emph{Probabilistic seasonal dengue forecasting in Vietnam: A modelling study using superensembles.}
PLOS Medicine 18(3): e1003542. \url{https://doi.org/10.1371/journal.pmed.1003542}. ლიცენზია: CC BY 4.0.

\noindent\textbf{ფაილი:} \texttt{colon\_gonzalez\_2021\_plosmed.pdf} (30 გვერდი; PDF-ის გვერდი = ჟურნალის გვერდი).

\tableofcontents

\section{სტატია ერთ აბზაცში}

ავტორებმა ვიეტნამისთვის \textbf{დენგეს პროგნოზის სისტემა} ააგეს, სახელად D-MOSS. სისტემა ყოველთვიურად
პროგნოზირებს, რამდენი შემთხვევა იქნება ქვეყნის 63 პროვინციიდან თითოეულში, \textbf{1-დან 6 თვემდე წინ}.
პროგნოზისთვის იყენებენ ძველ შემთხვევებს, სატელიტურ კლიმატის მონაცემებს და სეზონურ ამინდის პროგნოზებს.
ერთი მოდელის ნაცვლად რამდენიმე მოდელს აერთიანებენ (\q{superensemble}) და აჩვენებენ, რომ ეს ერთეულ მოდელებზე
უკეთესად მუშაობს. სისტემა ვიეტნამის ჯანდაცვის სამინისტროსთან და WHO-სთან ერთად შეიქმნა.

\begin{mnishvn}
შენი ნაშრომი \textbf{მოდელზე არ არის}, ის მოდელისთვის \textbf{მონაცემების მომზადებაზეა}. ამიტომ ამ სტატიიდან
მოდელის მათემატიკა არ გჭირდება. გჭირდება მხოლოდ ის, \textbf{როგორ მოამზადეს მონაცემები} მოდელამდე.
ეს მესამე მოდელირების სტატიაა (Suttirat-ისა და Lowe-ს შემდეგ), რომელიც ხელმძღვანელებმა მოითხოვეს:
ის აჩვენებს, რომ იგივე ნაბიჯები სხვა ქვეყანაშიც (ვიეტნამი) და სხვა ავტორებთანაც მეორდება.
\end{mnishvn}

\noindent\textbf{საინტერესო დეტალი:} ბოლო ავტორია \textbf{Rachel Lowe}, ბრაზილიის სტატიის (Lowe et al.\ 2021)
პირველი ავტორი. ეს ერთი კვლევითი ჯგუფია (London School of Hygiene \& Tropical Medicine).

\section{რა წაიკითხო}

მთლიანი სტატია 30 გვერდია, შენ მხოლოდ \textbf{დაახლოებით 6 გვერდი} გჭირდება:

\begin{center}
\begin{tabular}{@{}p{0.08\linewidth}p{0.42\linewidth}p{0.4\linewidth}@{}}
\toprule
გვ. & სექცია & რატომ \\
\midrule
1 & Data Availability Statement & ღიაა თუ არა მონაცემები \\
4--5 & Introduction, ბოლო აბზაცები & ორი ციტატა მოტივაციისთვის \\
6 & Dengue surveillance data & შემთხვევების წყარო \\
7 & Historical demographic \ldots; Historical Earth observation data & მოსახლეობა, კლიმატი, აგრეგაცია \\
8 & ზედა ნაწილი (საზღვრების ცვლილება) & \textbf{ყველაზე მნიშვნელოვანი} \\
9 & Model specification, ბოლო აბზაცი & lag-ები, ინსტრუმენტები \\
\bottomrule
\end{tabular}
\end{center}

\noindent\textbf{გამოტოვე:} ფორმულები (გვ.\ 8--11), Model selection, Superensemble, CRPS, Results (გვ.\ 12-დან).

\section{ნაწილ-ნაწილ: რას ამბობს სტატია}

\subsection{შემთხვევები (გვ.\ 6)}
\begin{citata}
\q{Monthly dengue cases were obtained from the Vietnamese Ministry of Health. Data were retrieved for the period
August 2002 to December 2019 at the province level (n = 63). Data comprised suspected and confirmed dengue cases,
although there was no indication as to how many cases fell within each category.}
\end{citata}
\noindent\textbf{მარტივად:} თვიური შემთხვევები 63 პროვინციისთვის, 2002--2019, საეჭვო და დადასტურებული ერთად.
გვ.\ 7-ზე ავტორები დასძენენ, რომ ამ პერიოდში შემთხვევის განსაზღვრება \textbf{არ შეცვლილა}.

\begin{mnishvn}
ეს ზუსტად ის საკითხია, რაც ტაილანდში ვნახეთ: OpenDengue-ში 2003 წლამდე „Total“ წერია, 2003-დან კი
„Suspected and confirmed“. ავტორები ცალკე ამოწმებენ, რომ განსაზღვრება დროში ერთნაირია, ანუ
\textbf{შედარებადობა} მონაცემების მომზადების ნაწილია. ეს ასაბუთებს შენს გადაწყვეტილებას (A20), რომ
სხვადასხვა განსაზღვრების წყაროები არ შეურიო.
\end{mnishvn}

\subsection{მოსახლეობა (გვ.\ 7)}
\begin{citata}
\q{Total population per province and per year were retrieved from the \ldots\ Gridded Population of the World
project version 4.11 \ldots\ 5 yearly for the period 2000 to 2020. Intervening years were generated using linear
interpolation.}
\end{citata}
\noindent\textbf{მარტივად:} მოსახლეობა მხოლოდ ყოველ 5 წელზე იყო ცნობილი, შუალედური წლები კი \textbf{წრფივად
შეავსეს}.

\begin{mnishvn}
ეს \textbf{შენი FR-17}-ია: შენი სისტემაც წლიურ მოსახლეობას წრფივად ინტერპოლირებს. მონახაზში წერია, რომ ეს
\q{eigene Entwurfsentscheidung}-ია, რადგან Suttirat სხვანაირად აკეთებდა. ახლა მას \textbf{მეცნიერული მხარდაჭერა}
აქვს: Colón-González-იც წრფივ ინტერპოლაციას იყენებს. ეს შეგიძლია ტექსტში დაამატო.
\end{mnishvn}

\subsection{კლიმატი: ბევრი წყარო (გვ.\ 7)}
ავტორები ერთ კვლევაში \textbf{ექვსზე მეტ} სხვადასხვა წყაროს იყენებენ, სხვადასხვა რეზოლუციით:
\begin{itemize}[nosep]
  \item ტემპერატურა (მინ., მაქს., საშუალო): MODIS-ის სატელიტი, 1 კმ;
  \item ნალექი: TRMM 2014 წლის აპრილამდე, \textbf{მერე GPM} (10 კმ);
  \item ტენიანობა: MODIS;
  \item ქარი: ERA5 2002--2011, \textbf{მერე NOAA CFS};
  \item El Niño-ს ინდექსი: NOAA.
\end{itemize}

\begin{mnishvn}
ორი რამ მნიშვნელოვანია:
\begin{enumerate}[nosep]
  \item \textbf{მოტივაციისთვის:} ერთი კვლევა, ექვსი წყარო, სხვადასხვა რეზოლუცია. ეს აჩვენებს, რამდენი შრომა
        სჭირდება მონაცემების შეკრებას. ზუსტად ამას ამარტივებს შენი სისტემა.
  \item \textbf{წყაროს შეცვლა დროში:} ნალექისა და ქარისთვის სერიის შუაში წყარო შეიცვალა. თუ ორი წყარო
        განსხვავებულად ზომავს, სერიაში „ნახტომი“ შეიძლება გაჩნდეს. ეს ზუსტად ის რისკია, რისთვისაც შენს სისტემაში
        \textbf{ორი წყაროს შედარება (FR-27)} არსებობს: ჯერ შეამოწმე, ემთხვევა თუ არა წყაროები, მერე გადააბი.
\end{enumerate}
\end{mnishvn}

\subsection{სივრცითი აგრეგაცია: რასტრიდან პროვინციამდე (გვ.\ 7--8)}
\begin{citata}
\q{Climate data were aggregated at the province level using population weighted averages for each month.}
\end{citata}

\begin{ganmarteba}
\textbf{მოსახლეობით შეწონილი საშუალო.} კლიმატის რუკა კვადრატებად (უჯრებად) არის დაყოფილი. პროვინციის
საშუალოს გამოსათვლელად ყველა უჯრა ერთნაირად კი არ ითვლება, არამედ იმის მიხედვით, \textbf{რამდენი ადამიანი
ცხოვრობს} თითოეულ უჯრაში. მაგალითი: პროვინციაში ორი უჯრაა, ქალაქი (30\,°C, 90\,000 ადამიანი) და მთა
(20\,°C, 10\,000 ადამიანი). უბრალო საშუალოა 25\,°C, მოსახლეობით შეწონილი კი
$(30\cdot 0{,}9 + 20\cdot 0{,}1) = 29\,$°C. ეს უფრო ზუსტად აღწერს, რა ტემპერატურაზე ცხოვრობს ადამიანების
უმეტესობა, ანუ სად ხდება დაავადების გადაცემა.
\end{ganmarteba}

\noindent წონებისთვის WorldPop-ის მონაცემები გამოიყენეს (100 მ). ეს მხოლოდ 2009, 2010, 2015 და 2020 წლებისთვის
იყო, ამიტომ შუალედური წლები ისევ \textbf{წრფივად ინტერპოლირეს}. ინსტრუმენტი: Python-ის მოდული
\texttt{rasterstats}.

\begin{mnishvn}
ახლა ორი ვარიანტი გაქვს ლიტერატურაში:
\begin{itemize}[nosep]
  \item \textbf{Lowe et al.\ (2021):} ფართობით შეწონილი საშუალო (\texttt{exactextractr}, R);
  \item \textbf{Colón-González et al.\ (2021):} მოსახლეობით შეწონილი საშუალო (\texttt{rasterstats}, Python).
\end{itemize}
ეს ზუსტად შენი \textbf{ვარიანტი B}-ია (ERA5-ის ბადე + შეწონილი საშუალო). დღევანდელმა ტესტმა აჩვენა, რომ
Open-Meteo-ს წერტილებით (ვარიანტი A) ეს ზედმეტად ძვირია, ამიტომ B მთავარი მიზანია. სისტემამ შეიძლება
მომხმარებელს \textbf{არჩევანი} მისცეს: ფართობით თუ მოსახლეობით. ეს ახალი მოთხოვნის კანდიდატია.
\end{mnishvn}

\subsection{საზღვრების ცვლილება (გვ.\ 8): ყველაზე მნიშვნელოვანი}
\begin{citata}
\q{There were multiple changes in the administrative boundaries of the provinces over the study period. The last
change took place in 2008. Consequently, we froze the administrative boundaries of the country at their 2008 level.
Dengue data were allocated to each of the 2008 level administrative boundaries using district-level records.}
\end{citata}
\noindent\textbf{მარტივად:} ვიეტნამში კვლევის პერიოდში პროვინციების საზღვრები რამდენჯერმე შეიცვალა. ავტორებმა
\textbf{ერთი წლის საზღვრები აირჩიეს} (2008) და ყველა წლის შემთხვევები ამ საზღვრებზე გადაანაწილეს, რაიონების
ჩანაწერების მიხედვით (ეს ჯანდაცვის სამინისტრომ გააკეთა).

\begin{mnishvn}
ეს \textbf{შენი FR-28-ის ბირთვია}. ახლა ერთი და იგივე პრობლემა სამ ადგილას გაქვს:
\begin{enumerate}[nosep]
  \item \textbf{Suttirat (ტაილანდი):} ახალი პროვინცია (Bueng Kan, 2011) ძველს \textbf{შეუერთეს}. შენი სისტემა უკვე
        აკეთებს ამას (დავალება H, 18\,240 პროვინცია-თვე).
  \item \textbf{Colón-González (ვიეტნამი):} საზღვრები ერთ წელზე \textbf{„გაყინეს“} (2008).
  \item \textbf{HDX (ვიეტნამი, 2025):} 63 პროვინცია 34-ად გაერთიანდა, ამიტომ დღევანდელი რუკა ძველ მონაცემებს
        არ ემთხვევა.
\end{enumerate}
\textbf{დასკვნა თეზისისთვის:} „დროში ცვალებადი ადმინისტრაციული ერთეულების ჰარმონიზაცია“ ერთი სტატიის
თავისებურება კი არა, \textbf{ტიპური ნაბიჯია}. ეს Traceability-Matrix-ში FR-28-ის მეორე წყაროდ ჩაიწერება.
ორი სტრატეგია არსებობს (შეერთება ან ერთ წელზე გაყინვა), და სისტემამ ორივე შეიძლება შესთავაზოს.
\end{mnishvn}

\subsection{lag-ები (გვ.\ 4, 9)}
\begin{citata}
\q{Most of the effects of climate on dengue incidence are delayed by 1 to 2 months \ldots} (გვ.\ 4)
\end{citata}
\noindent გვ.\ 9-ზე: ტემპერატურისთვის, ნალექისთვის, ტენიანობისა და DTR-ისთვის (დღე-ღამის ტემპერატურის სხვაობა)
\textbf{3-თვიანი მოძრავი საშუალო} გამოიყენეს, lag 0--2 თვე. El Niño-სთვის lag 0--3. ქარისთვის lag არ გამოიყენეს.
ასევე \textbf{წინა თვის შემთხვევები} მოდელის შესავალია.

\begin{mnishvn}
შენს სისტემას lag-სვეტები აქვს (FR-23), მაგრამ \textbf{მოძრავი საშუალო} (მაგ.\ lag 0--2-ის საშუალო) და
\textbf{lag-იანი შემთხვევები} (წინა თვის შემთხვევები) არა. ორივე ადვილი დასამატებელია და კარგი კანდიდატია
შემდეგი ნაბიჯისთვის, ან Future Work-ისთვის.
\end{mnishvn}

\subsection{ინსტრუმენტები და მონაცემების ხელმისაწვდომობა (გვ.\ 1, 7, 9)}
\begin{itemize}[nosep]
  \item Python 3.6 (\texttt{rasterstats}) და R 3.6.1 + INLA. ესე იგი \textbf{ორი პროგრამირების ენა} ერთ კვლევაში;
  \item კოდი GitHub-ზეა (\url{https://github.com/FelipeJColon/paper_dengue_superensemble});
  \item დენგეს მონაცემები \textbf{ღია არ არის}: მხოლოდ მოთხოვნით, ვიეტნამის სამინისტროს კონკრეტული პირისგან.
\end{itemize}

\begin{mnishvn}
\q{მონაცემები მოთხოვნით} ნიშნავს, რომ სხვა მკვლევარს იგივე კვლევის გამეორება უჭირს. შენი სისტემა
\textbf{ღია წყაროებს} (OpenDengue და სხვ.) იყენებს და ყოველ ნაბიჯს აღრიცხავს. ეს \textbf{განმეორებადობის}
(reproducibility) არგუმენტია, რომელიც შენს კვლევით კითხვაშიც არის.
\end{mnishvn}

\section{ორი ციტატა მოტივაციისთვის}
\begin{citata}
\q{Some of these challenges include the lack of multi-decadal health data sets with which to train and validate
seasonal climate-driven early warning systems, \textbf{the common mismatch of scales between climate data outputs and
data used for decision-making}, and a general lack of consensus as to how to communicate uncertainties to users.}
(გვ.\ 4)
\end{citata}
\noindent \textbf{რას ნიშნავს:} კლიმატის მონაცემები ბადეზეა, ჯანდაცვის მონაცემები კი პროვინციებზე. მათ
„მასშტაბები“ არ ემთხვევა. ზუსტად ამას აერთიანებს შენი სისტემა.

\begin{citata}
\q{\ldots difficulties related to the integration of these systems into existing public health procedures and
\textbf{occasional lack of technical expertise}.} (გვ.\ 5)
\end{citata}
\noindent \textbf{რას ნიშნავს:} ზოგჯერ ტექნიკური ცოდნა აკლიათ. ეს Boylan-ისა და Rajamani-ს გვერდით უნარების
ნაკლებობის \textbf{მესამე, ფრთხილი} წყაროა. ფრთხილი იმიტომ, რომ ავტორები \q{occasional}-ს წერენ, ანუ „ზოგჯერ“.
ტექსტში არ გაამძაფრო.

\section{შენი ცხრილისთვის (Schritte / Traceability-Matrix)}

\begin{longtable}{@{}p{0.27\linewidth}p{0.33\linewidth}p{0.06\linewidth}p{0.26\linewidth}@{}}
\toprule
ნაბიჯი & Colón-González 2021 & გვ. & შენი სისტემა \\
\midrule
\endhead
თვიური შემთხვევები & ჯანდაცვის სამინისტრო, 63 პროვინცია, 2002--2019 & 6 & \ok{} (FR-28, ტაილანდი) \\
შემთხვევის განსაზღვრება დროში & არ შეცვლილა (შემოწმებულია) & 6--7 & \ok{} FR-26/27, A20 \\
საზღვრების ჰარმონიზაცია & საზღვრები 2008-ზე გაყინული & 8 & \ok{} ტაილანდი, \warn{} ზოგადად \\
კლიმატის რამდენიმე წყარო & TRMM$\rightarrow$GPM, ERA5$\rightarrow$CFS & 7 & \warn{} შედარება \ok{}, გადაბმა 
o{} \\
რასტრი $\rightarrow$ პროვინცია & მოსახლეობით შეწონილი (\texttt{rasterstats}) & 7 & 
o{} (FR-28, ვარიანტი B) \\
მოსახლეობის შუალედური წლები & წრფივი ინტერპოლაცია & 7 & \ok{} FR-17 \\
lag-ები & 3-თვიანი საშუალო lag 0--2, ENSO 0--3 & 9 & \ok{} სვეტები, 
o{} მოძრავი საშუალო \\
ინსტრუმენტები & Python + R/INLA & 7, 9 & -- \\
\bottomrule
\end{longtable}

\section{შენი დავალება}
\begin{enumerate}[nosep]
  \item წაიკითხე ზემოთ ჩამოთვლილი 6 გვერდი და ციტატები PDF-ში გადაამოწმე (Ctrl+F).
  \item შეავსე კითხვის ფურცელი (\texttt{kitxvis\_furceli\_colon\_gonzalez\_2021.md}) \textbf{შენი სიტყვებით}.
  \item ბოლოს უპასუხე: რას ამატებს ეს სტატია ჩემს ნაშრომს? (მინიშნება: საზღვრები, მოსახლეობით შეწონვა,
        წყაროს ცვლილება, FR-17-ის მხარდაჭერა.)
  \item მომწერე: \q{შეამოწმე Colón-González}.
\end{enumerate}

\end{document}
